\phantomsection
\color{uniblue}\section*{\centering{\large{ПЕРЕЧЕНЬ СОКРАЩЕНИЙ}\color{white!0}<.>}}
\addcontentsline{toc}{section}{Перечень сокращений}
\color{black}
\begin{longtable}{>{\raggedright\arraybackslash}m{2cm}>{\raggedright\arraybackslash}m{0.5cm}>{\raggedright\arraybackslash}m{20cm}}
\endfirsthead\endhead\endfoot\endlastfoot
DSP & -- & Digital Signal Processing; \\
GOOSE & -- & Generic Object-Oriented Substation Event; \\
IRF & -- & Internal Relay Fault; \\
PPS & -- & Pulse Per Second; \\
SPI & -- & Serial Peripheral Interface; \\
АВР & -- & автоматическое включение резерва; \\
АПВ & -- & автоматическое повторное включение; \\
АСУ~ТП & -- & автоматизированная система управления технологическими процессами; \\
АЦП & -- & аналого-цифровой преобразователь; \\
БИ & -- & блок испытательный; \\
БЛЗШ & -- & блокировка логической защиты шин; \\
БНН & -- & блокировка при неисправности цепей напряжения; \\
БНТ & -- & блокировка при бросках намагничивающего тока; \\
ВН & -- & высшее напряжение; \\
ГЗ & -- & газовая защита; \\
ДО & -- & дочерние общества; \\
ДТм & -- & датчик температуры масла; \\
ДТо & -- & датчик температуры обмотки; \\
ДУм & -- & датчик уровня масла; \\
ЗАВР & -- & запрет автоматического включения резерва; \\
ЗАПВ & -- & запрет автоматического повторного включения; \\
ЗДЗ & -- & защита от дуговых замыканий; \\
ЗОП & -- & защита от обрыва провода; \\
ЗП & -- & защита от перегрузки; \\
ИО & -- & измерительный орган; \\
ИЧМ & -- & интерфейс человек-машина; \\
ИЭУ & -- & интеллектуальное электронное устройство; \\
КА & -- & коммутационный аппарат; \\
КЗ & -- & короткое замыкание; \\
КИ & -- & контроль изоляции; \\
КП & -- & контроллер присоединения; \\
КПОН & -- & комбинированный пусковой орган напряжения; \\
КПруж & -- & контроль пружины привода выключателя; \\
КРВ & -- & контроль ресурса выключателя; \\
КС & -- & контроль синхронизма / контрольная сумма; \\
КСВ & -- & контроль силового выключателя; \\
КЦВ & -- & контроль цепей выключателя; \\
КЦН & -- & контроль цепей напряжения; \\
ЛЗТ & -- & логическая защита трансформатора; \\
ЛЗШ & -- & логическая защита шин; \\
ЛО & -- & логика отключения; \\
МРВ & -- & механический ресурс выключателя; \\
МТЗ & -- & максимальная токовая защита; \\
НКУ & -- & низковольтное комплектное устройство; \\
НН & -- & низшее напряжение; \\
НП & -- & нулевая последовательность; \\
ОВ & -- & оперативный вывод; \\
ОЗУ & -- & оперативное запоминающее устройство; \\
ОК & -- & отсечной клапан; \\
ОП & -- & обратная последовательность; \\
ОТ & -- & оперативный ток; \\
ПДС & -- & преобразователь дискретных сигналов; \\
ПЗУ & -- & постоянное запоминающее устройство; \\
ПК & -- & предохранительный клапан; \\
ПО & -- & пусковой орган / программное обеспечение; \\
ПП & -- & прямая последовательность; \\
ПС & -- & предупредительная сигнализация; \\
ПУ & -- & пульт управления; \\
РАС & -- & регистратор аварийных событий; \\
РД & -- & реле давления; \\
РЗА & -- & релейная защита и автоматика; \\
РПН & -- & устройство регулирования напряжения под нагрузкой; \\
РФК & -- & реле фиксации команд; \\
РЭ & -- & руководство по эксплуатации; \\
СВ & -- & секционный выключатель; \\
СС & -- & сборка сигналов; \\
ТЗ & -- & технологические защиты; \\
ТК & -- & токовый контроль; \\
ТН & -- & трансформатор напряжения; \\
ТО & -- & токовая отсечка; \\
ТС & -- & технологическая сигнализация; \\
ТТ & -- & трансформатор тока; \\
ТУ & -- & телеуправление/телеускорение; \\
УВ & -- & управление выключателем / управляющее воздействие; \\
УРОВ & -- & устройство резервирования при отказе выключателя; \\
ФСУ & -- & функциональная схема устройства; \\
ЦН & -- & цепи напряжения; \\
ЦП & -- & центральный процессор; \\
ШЭТ & -- & шкаф электротехнический типовой; \\
ЭМ & -- & электромагнит; \\
ЭМВ & -- & электромагнит включения; \\
ЭМО & -- & электромагнит отключения; \\
ЭМУ & -- & электромагнит управления. \\
\end{longtable}